\begin{theorem}[The maximal regular-locus
\texorpdfstring{\(\mathbb R^4\)}{R4} quotient and its compactification
obstruction]
\label{cusp-spec-globalization-obstruction}
\mbox{\Verified}\ 
Let
\[
 \Lambda=\mathbb Z\langle\widehat\gamma,\widehat u,
                   \widehat w,\widehat\delta\rangle,\qquad
 K_\Lambda=\mathbb Z\langle\widehat w,\widehat\delta\rangle,
\]
and define
\[
 \varepsilon_\gamma
 (a\widehat\gamma+b\widehat u+c\widehat w+d\widehat\delta)=a,\qquad
 L_\Lambda=\ker\varepsilon_\gamma
 =\mathbb Z\langle\widehat u,\widehat w,\widehat\delta\rangle.
\tag{G.1}
\]
Then the following statements hold.

\smallskip
\noindent
\emph{(i) Regular group quotient.}
Over \(B^\circ\), the period family \(J=X^\circ=f^{-1}(B^\circ)\) is a
holomorphic family of complex Lie groups, not merely an affine torus
bundle.  The formula
\begin{equation}
 \chi_\gamma\bigl([z,\Pi(z)x]\bigr)
 =\bigl([z],\varepsilon_\gamma(x)\bmod\mathbb Z\bigr)
 \label{cusp-spec-global-circle-map}
\end{equation}
defines a global smooth real-analytic fibrewise group epimorphism
\begin{equation}
 \chi_\gamma:J\longrightarrow B^\circ\times(\mathbb R/\mathbb Z).
 \label{cusp-spec-global-circle-quotient}
\end{equation}
Write
\begin{equation}
 \theta_\gamma:=\operatorname{pr}_{\mathbb R/\mathbb Z}\circ\chi_\gamma:
 J\longrightarrow\mathbb R/\mathbb Z
 \label{cusp-spec-global-circle-coordinate}
\end{equation}
for its circle coordinate.  Thus
\(\chi_\gamma=(f,\theta_\gamma)\); the two symbols have different
codomains and will not be interchanged below.
Its kernel is the connected real three-torus family
\begin{equation}
 \mathcal H:=
 \frac{\mathfrak h_z^\circ\times
       ((L_\Lambda\otimes\mathbb R)/L_\Lambda)}{\Delta},
 \label{cusp-spec-global-kernel-family}
\end{equation}
where \(\Delta=\pi_1(B^\circ)\) acts diagonally through the deck action
on \(\mathfrak h_z^\circ\) and the restricted monodromy on
\(L_\Lambda\).  Thus every fibre of \(\mathcal H\) is
\((L_\Lambda\otimes\mathbb R)/L_\Lambda\); the family need not be a
product.
The quotient is real and is not a holomorphic quotient of complex tori.

\smallskip
\noindent
\emph{(ii) Exact radial quotient metric.}
On the positive radial cusp, the restriction of
\cref{cusp-spec-global-circle-map} to
\(F_T=\mathbb R^4/P_T\mathbb Z^4\) is
\begin{equation}
 \chi_{\gamma,T}([y])
 =\left\langle P_T^{-\mathsf T}e_\gamma,y\right\rangle
   \bmod\mathbb Z,\qquad e_\gamma=(1,0,0,0)^{\mathsf T}.
 \label{cusp-spec-radial-circle-character}
\end{equation}
For the Riemannian quotient of the Euclidean metric, the target circle has
circumference
\begin{equation}
 \ell_T
 =\frac1{\lVert P_T^{-\mathsf T}e_\gamma\rVert}
 =\frac{D_T}{\sqrt{L_T^2+m_T^2}}
 =T+H-B+O(T^{-1}).
 \label{cusp-spec-global-circle-circumference}
\end{equation}
In particular, \(\ell_T\to\infty\).

\smallskip
\noindent
\emph{(iii) Four copies and minimality.}
The source
\begin{equation}
 J^{[4]}=J\times_{B^\circ}J\times_{B^\circ}J\times_{B^\circ}J
 \label{cusp-spec-fourfold-fibre-product}
\end{equation}
is the actual fourfold fibre product over all of \(B^\circ\).  The product
map
\begin{equation}
 \chi_\gamma^{[4]}=(\chi_\gamma,\chi_\gamma,\chi_\gamma,\chi_\gamma):
 J^{[4]}\longrightarrow B^\circ\times(\mathbb R/\mathbb Z)^4
 \label{cusp-spec-fourfold-circle-map}
\end{equation}
has kernel
\(\mathcal H^{[4]}=\mathcal H\times_{B^\circ}\cdots
\times_{B^\circ}\mathcal H\), whose fibres are
\(\mathbb T^{12}\).  Along the synchronized positive
radial cusp, its quotient fibre, with product quotient metric, is
\begin{equation}
 G_T=(\mathbb R/\ell_T\mathbb Z)^4,\qquad
 \operatorname{sys}(G_T)=\ell_T.
 \label{cusp-spec-four-circle-fibre}
\end{equation}
Consequently
\begin{equation}
 (G_T,0)\longrightarrow(\mathbb R^4,0)
 \label{cusp-spec-four-circle-limit}
\end{equation}
in pointed local-isometric, pointed smooth, and pointed
Gromov--Hausdorff topology.  If \(\ell_T>4R\), the radius-\(R\) ball
about zero is exactly Euclidean.

Equivalently, the external Cartesian product gives a quotient
\[
 J^4\longrightarrow(B^\circ)^4\times(\mathbb R/\mathbb Z)^4,
\]
and \(J^{[4]}\) is its pullback along the diagonal
\(B^\circ\to(B^\circ)^4\).  At four independent cusp parameters its target
lattice is
\(\operatorname{diag}(\ell_{T_1},\ldots,\ell_{T_4})\mathbb Z^4\), so the
same pointed limit holds whenever
\(\min_i\ell_{T_i}\to\infty\).

Among connected torus quotients obtained from \(r\) copies by a primitive
full-monodromy-invariant kernel and whose quotient systole tends to infinity
on the radial cusp, the quotient rank is at most \(r\).  Thus four copies
are minimal for a four-dimensional all-expanding quotient, and
\cref{cusp-spec-fourfold-circle-map} attains the bound.

The cusp-local two-copy map
\begin{equation}
 \pi_T\times\pi_T:F_T\times F_T\longrightarrow Q_T\times Q_T
 \label{cusp-spec-doubled-local-quotient}
\end{equation}
from \cref{cusp-spec-native-quotient} also has pointed limit
\(\mathbb R^4\).  It is not a quotient over the full \(B^\circ\), however,
and no finite base cover makes it one.

\smallskip
\noindent
\emph{(iv) Finite fillings.}
At \(p_j\), use the notation of
\cref{prop:candidate-finite-fillings} and put
\[
 (m_1,v_1,\kappa_1)
 =\bigl(3,\widehat\gamma+2\widehat u-4\widehat w,1\bigr),
\]
\[
 (m_2,v_2,\kappa_2)
 =\bigl(4,-\widehat\gamma-3\widehat u+3\widehat w,-1\bigr),\qquad
 \kappa_j=\varepsilon_\gamma(v_j).
\tag{G.2}
\]
In the un-conjugated flat coordinates, the finite filling generator is
\begin{equation}
 \Phi_j(s,x)=
 \left(e^{-2\pi i/m_j}s,A_jx+\frac{v_j}{m_j}\right).
 \label{cusp-spec-finite-affine-action}
\end{equation}
The raw character is not invariant:
\begin{equation}
 \varepsilon_\gamma\left(A_jx+\frac{v_j}{m_j}\right)
 =\varepsilon_\gamma(x)+\frac{\kappa_j}{m_j}\pmod{\mathbb Z}.
 \label{cusp-spec-finite-character-shift}
\end{equation}
It instead descends equivariantly to the Seifert circle family
\begin{equation}
 C_j=
 \frac{\Delta_j\times(\mathbb R/\mathbb Z)}
 {(s,\theta)\sim
  (e^{-2\pi i/m_j}s,\theta+\kappa_j/m_j)}.
 \label{cusp-spec-finite-Seifert-target}
\end{equation}
The total space \(C_j\) is a smooth solid torus, the map
\(C_j\to D_j\) has a multiple central circle of order \(m_j\), and the
induced map \(N_j\to C_j\) is a smooth proper \(\mathbb T^3\)-bundle.
The target translation in \cref{cusp-spec-finite-Seifert-target} is
exactly compatible with the logarithmic gluing map \(h_j\).

For four synchronized copies, the smooth simultaneous-root extension is
\begin{equation}
 \widetilde N_j^{[4]}
 =\frac{(J'_j)^{\times_{\Delta_j}4}}
 {\langle(\widetilde g_j^{\log})^{\times4}\rangle},
 \label{cusp-spec-finite-fourfold-normalization}
\end{equation}
with target
\begin{equation}
 C_j^{[4]}=
 \frac{\Delta_j\times(\mathbb R/\mathbb Z)^4}
 {(s,\theta_1,\ldots,\theta_4)\sim
  (e^{-2\pi i/m_j}s,
   \theta_1+\kappa_j/m_j,\ldots,\theta_4+\kappa_j/m_j)}.
 \label{cusp-spec-finite-fourfold-target}
\end{equation}
The map \(\widetilde N_j^{[4]}\to C_j^{[4]}\) is a smooth proper
\(\mathbb T^{12}\)-bundle.  There is a canonical finite map
\begin{equation}
 \nu_j:\widetilde N_j^{[4]}
 \longrightarrow N_j\times_{D_j}N_j\times_{D_j}N_j\times_{D_j}N_j
 \label{cusp-spec-finite-normalization-map}
\end{equation}
which is an isomorphism over \(D_j^*\), is the normalization of the
reduced ordinary filled fibre product, and has degree \(m_j^3\) on the
reduced central fibre.  These two degrees are \(27\) and \(64\).  Gluing
these two local normalizations to \(J^{[4]}\) by
\((h_j)^{\times4}\), and gluing their targets to the product
\(B^\circ\times(\mathbb R/\mathbb Z)^4\) by
\((h_j^{S^1})^{\times4}\), gives the corresponding normalization and
Seifert target over \(B\setminus\{p_0\}\).

\smallskip
\noindent
\emph{(v) Cusp nonextension.}
Neither the literal map \(\chi_\gamma\) nor
\(\chi_\gamma^{[4]}\) extends continuously through the toric cusp fibre
\(W\).  Thus the four-copy quotient is global over \(B^\circ\), and its
finite-filling normalization is global over \(B\setminus\{p_0\}\), but
the displayed map \(\chi_\gamma^{[4]}\) has no continuous extension to
the compact fourfold fibre product or to \(X^4\) with the stated
regular-fibre restriction, and it does not factor through
\(X^{\wedge4}\cong S^{24}\).  The cohomology class of
\(\chi_\gamma\) does extend over
\(N_0\); this weaker homotopy statement does not extend the displayed
quotient map or its metric.

\smallskip
\noindent
\emph{(vi) Products, smashes, and sphere constructions.}
If the recognition \(X\cong_{\mathrm{diff}}S^6\) is used, then
\begin{equation}
 X^2\cong_{\mathrm{diff}}S^6\times S^6\not\simeq S^{12},\qquad
 H_6(X^2;\mathbb Z)\cong\mathbb Z^2,
 \label{cusp-spec-product-not-sphere}
\end{equation}
and
\begin{equation}
 \begin{aligned}
 X^4&\cong_{\mathrm{diff}}(S^6)^4\not\simeq S^{24},\\
 H_6(X^4;\mathbb Z)&\cong\mathbb Z^4,&
 H_{12}(X^4;\mathbb Z)&\cong\mathbb Z^6,&
 H_{18}(X^4;\mathbb Z)&\cong\mathbb Z^4.
 \end{aligned}
 \label{cusp-spec-four-product-not-sphere}
\end{equation}
The synchronized regular fourfold \(J^{[4]}\) has real dimension \(18\),
not \(24\).  Choosing a base point gives based homeomorphisms
\begin{equation}
X\wedge X\cong S^{12},\qquad X^{\wedge4}\cong S^{24},
\label{cusp-spec-smash-spheres}
\end{equation}
but these smash spheres do not carry either \emph{literal fixed-target}
quotient above.  The later cone constructions
\cref{cusp-cone-S12-route,cusp-cone-S24-local-route,cusp-cone-S24-route}
change the codomain by retaining the radial parameter and collapsing the
obstructing boundary cross-section; they therefore do not contradict this
nonextension theorem.  Moreover,
neither \(S^{12}\) nor \(S^{24}\), with any differentiable structure,
admits an almost-complex structure
\cite[\S15, Proposition~15.1 and the following remark,
p.~434]{BorelSerre1953}; hence neither can be promoted to a compact
complex manifold.  This is an obstruction to a proposed complex
\(S^{12}\)- or \(S^{24}\)-globalization, not to the real quotient or its
pointed metric limit.  Precisely,
\(\pi_T\times\pi_T\) is a map on \(F_T^2\), and it fails to factor
through \(F_T\wedge F_T\); \(\chi_{\gamma,T}^{[4]}\) is a map on
\(F_T^4\), and it fails to factor through \(F_T^{\wedge4}\).
Neither map is a map on \(X^2\) or \(X^4\), and no extension with the
stated fibre restriction exists.

There is one exact but different sphere construction.  If
\(C_T=\mathbb R/\ell_T\mathbb Z\) is based at zero, postcomposition with
the target fat-wedge collapse gives
\begin{equation}
 F_T^4\xrightarrow{\chi_{\gamma,T}^{[4]}}C_T^4
 \longrightarrow C_T^{\wedge4}\cong S^4.
 \label{cusp-spec-target-smash-map}
\end{equation}
This map does factor through \(F_T^{\wedge4}\).  With the quotient length
metric and base point
\([\ell_T/2,\ell_T/2,\ell_T/2,\ell_T/2]\), its pointed limit is
\(\mathbb R^4\), with every radius-\(R\) ball exactly Euclidean once
\(\ell_T>4R\).  The metric is singular at the collapsed fat wedge, and
the map still does not extend to \(X^4\) or \(X^{\wedge4}\).
\end{theorem}

\begin{proof}
\medskip\noindent\emph{Invariant lattices and the regular group law.}\par
The matrices in \cref{prop:monodromy-certificate} give
\[
 (M_0-I)(a,b,c,d)^{\mathsf T}=(0,0,b,-a)^{\mathsf T},
\]
so
\begin{equation}
 K_\Lambda=\ker(M_0-I)=\operatorname{im}(M_0-I).
 \label{cusp-spec-cusp-kernel}
\end{equation}
They also give
\begin{equation}
 \begin{aligned}
 A_1\widehat u&=-\widehat w+\widehat\delta,&
 A_1\widehat w&=\widehat u-\widehat w,&
 A_1\widehat\delta&=\widehat\delta,\\
 A_2\widehat u&=\widehat w,&
 A_2\widehat w&=-\widehat u+\widehat\delta,&
 A_2\widehat\delta&=\widehat\delta.
 \end{aligned}
 \label{cusp-spec-invariant-flag-calculation}
\end{equation}
Thus \(L_\Lambda\) is invariant.  It is the smallest invariant lattice
containing \(K_\Lambda\), because
\(\widehat u=A_1\widehat w+\widehat w\).  Both lattices are primitive.

For a row covector \((a,b,c,d)\), invariance under \(A_1\) gives
\(b=2c\) and \(d=3c\), whereas invariance under \(A_2\) gives
\(b=c\) and \(d=2c\).  Hence \(b=c=d=0\), and
\begin{equation}
 \operatorname{Hom}(\Lambda,\mathbb Z)^\Delta
 =\mathbb Z\varepsilon_\gamma.
 \label{cusp-spec-invariant-character-lattice}
\end{equation}

The family over the cover is
\((\mathfrak h_z^\circ\times\mathbb C^2)/\Lambda\), and the action used
to form \(J\) is
\[
 [z,\zeta]\longmapsto[gz,R_g(z)\zeta].
\]
It is complex linear in every fibre and fixes the zero section.
Fibrewise addition, inversion, and zero are therefore equivariant and
descend to \(J\).  In the flat real-analytic marking
\[
 \Xi(z,x+\Lambda)=[z,\Pi(z)x],
\]
equation (2.3d) makes the action
\[
 (z,x+\Lambda)\longmapsto(gz,A_gx+\Lambda).
\]
Equation \cref{cusp-spec-invariant-character-lattice} proves that
\cref{cusp-spec-global-circle-map} is well defined.  Primitivity of
\(\varepsilon_\gamma\) proves surjectivity and connectedness of the
kernel.  The induced action on \(\Lambda/L_\Lambda\simeq\mathbb Z\) is
trivial, so the target is the product circle bundle.  Since the kernel has
odd real dimension, this is not a holomorphic quotient.

\medskip\noindent\emph{The quotient metric and the four-copy limit.}\par
The first column of
\cref{cusp-spec-full-inverse-transpose} is
\begin{equation}
 P_T^{-\mathsf T}e_\gamma
 =\left(0,\frac{m_T}{D_T},0,-\frac{L_T}{D_T}\right)^{\mathsf T}.
 \label{cusp-spec-circle-dual-vector}
\end{equation}
The quotient metric on \(\mathbb R/\mathbb Z\) is therefore
\(\ell_T^2d\theta^2\).  Equivalently,
\[
 P_T(L_\Lambda\otimes\mathbb R)
 =\operatorname{span}_{\mathbb R}\{e_1,e_3,v_{u,T}\}
\]
has unit normal
\[
 \nu_T=\frac{(0,m_T,0,-L_T)}{\sqrt{L_T^2+m_T^2}},
\]
and
\[
 \operatorname{dist}
 \bigl(v_{\gamma,T},P_T(L_\Lambda\otimes\mathbb R)\bigr)
 =\left|\langle v_{\gamma,T},\nu_T\rangle\right|
 =\frac{L_Tq_T+6m_T^2}{\sqrt{L_T^2+m_T^2}}
 =\ell_T.
\]
The expansions in
\cref{cusp-spec-radial-notation,cusp-spec-holomorphic-errors} give
\[
 D_T=T^2+(2H-B)T+O(1),\qquad
 \sqrt{L_T^2+m_T^2}=T+H+O(T^{-1}),
\]
which proves \cref{cusp-spec-global-circle-circumference}.

After rescaling the circle coordinate by \(\ell_T\), the product target is
\(\mathbb R^4/\ell_T\mathbb Z^4\).  If
\(x,y\in B_R(0)\) and \(0\ne\lambda\in\ell_T\mathbb Z^4\), then
\[
 \lVert x-y+\lambda\rVert
 \ge\lVert\lambda\rVert-\lVert x-y\rVert
 >\ell_T-2R>2R\ge\lVert x-y\rVert
\]
when \(\ell_T>4R\).  Hence the covering projection is an exact isometry
on \(B_R(0)\), proving all three convergences in
\cref{cusp-spec-four-circle-limit}.  The quotient map is global over
\(B^\circ\); only the metric limit is evaluated along the radial cusp.
The separately chosen Euclidean fibre metric is not claimed to be a
canonical global metric away from that cusp.

\medskip\noindent\emph{Minimality and the finite-base-cover obstruction.}\par
Let \(N\subset\Lambda^{\oplus r}\) be a primitive invariant kernel whose
connected quotient has radial systole tending to infinity.  Every element
of \(K_\Lambda^{\oplus r}\) has uniformly bounded physical length, since
\[
 P_T\widehat w=e_1,\qquad P_T\widehat\delta=e_3.
\]
If such an element survived in the quotient lattice, its quotient norm
would be no greater than its source norm and would give a nonzero uniformly
bounded period.  Therefore \(K_\Lambda^{\oplus r}\subset N\).
Invariance and
\cref{cusp-spec-invariant-flag-calculation} then force
\(L_\Lambda^{\oplus r}\subset N\), and
\[
 \operatorname{rank}(\Lambda^{\oplus r}/N)\le4r-3r=r.
\]
This proves the lower bound \(r\ge4\), while
\(N=L_\Lambda^{\oplus4}\) proves sharpness.

For the finite-cover assertion, put
\begin{equation}
 H_*=A_2A_1
 =\begin{pmatrix}
 1&0&0&0\\
 6&1&1&0\\
 0&0&1&0\\
 -5&0&-1&1
 \end{pmatrix}=I+N_*.
 \label{cusp-spec-no-finite-cover-matrix}
\end{equation}
Direct multiplication gives \(N_*^2=0\), and hence
\begin{equation}
 H_*^n\widehat w
 =\widehat w+n(\widehat u-\widehat\delta),\qquad
 H_*^n\widehat\delta=\widehat\delta.
 \label{cusp-spec-no-finite-cover-orbit}
\end{equation}
Thus \(H_*^nK_\Lambda\ne K_\Lambda\) for every \(n\ne0\).
Every finite-index subgroup of the base fundamental group contains a
positive power of the loop with monodromy \(H_*\).  Such a subgroup cannot
preserve \(K_\Lambda\), proving that no finite base cover globalizes the
native two-torus quotient.

\medskip\noindent\emph{The affine finite fillings and the actual overlap
maps.}\par
Equation \cref{cusp-spec-invariant-character-lattice} gives
\(\varepsilon_\gamma A_j=\varepsilon_\gamma\), so
\cref{cusp-spec-finite-character-shift} follows.  Define
\[
 \overline\Phi_j(s,\theta)
 =\left(e^{-2\pi i/m_j}s,\theta+\frac{\kappa_j}{m_j}\right).
\]
Then \((s,x)\mapsto(s,\varepsilon_\gamma(x))\) is
\((\Phi_j,\overline\Phi_j)\)-equivariant; denote the descended map by
\(\overline\chi_j:N_j\to C_j\).  A nonidentity power which fixes
the base must lie over \(s=0\), where it translates the circle by
\(k\kappa_j/m_j\).  Since \(\gcd(\kappa_j,m_j)=1\), this is nonintegral
for \(0<k<m_j\).  The target action is free; the source action is free by
\cref{prop:candidate-finite-fillings}.

For an explicit solid-torus chart, put
\(\lambda_j=e^{2\pi i/m_j}\), \(u=e^{2\pi i\theta}\), and choose
\(r_j\in\mathbb Z\) with \(r_j\kappa_j\equiv1\pmod {m_j}\).  Then
\[
 w=su^{r_j},\qquad U=u^{m_j}
\]
are invariant and identify \(C_j\) real-analytically with a solid torus.
In these coordinates
\[
 t_j=s^{m_j}=w^{m_j}U^{-r_j},
\]
which proves the central multiplicity.  Equivariance and the primitive
kernel give the proper \(\mathbb T^3\)-bundle \(N_j\to C_j\).

It remains to check the actual logarithmic gluing.  Let
\(\Pi_{\mathbb R}(z):\Lambda\otimes\mathbb R\to\mathbb C^2\) be the
underlying real isomorphism.  On a logarithm branch \(\ell\), define
\(b_j^\ell(z)\) uniquely by
\begin{equation}
 \Pi_{\mathbb R}(z)b_j^\ell(z)
 =\frac{\ell(z)}{2\pi i}\Pi(z)v_j,\qquad
 \beta_j(z)=\varepsilon_\gamma(b_j^\ell(z))\bmod\mathbb Z.
 \label{cusp-spec-finite-target-conjugacy}
\end{equation}
Changing \(\ell\) by \(2\pi in\) adds \(nv_j\) to \(b_j^\ell\), hence
adds \(n\kappa_j\in\mathbb Z\) to its
\(\varepsilon_\gamma\)-value.  Thus \(\beta_j\) is branch independent.
In flat coordinates the source conjugacy (2.9e) is translation by
\(b_j^\ell\).  Applying \(\varepsilon_\gamma\) to that conjugacy gives
\begin{equation}
 \beta_j(g_jz)=\beta_j(z)-\frac{\kappa_j}{m_j}\pmod{\mathbb Z}.
 \label{cusp-spec-finite-target-cocycle}
\end{equation}
Therefore
\[
 h_j^{S^1}(z,\theta)=(z,\theta+\beta_j(z))
\]
conjugates \(\overline\Phi_j\) to the linear target action
\((z,\theta)\mapsto(g_jz,\theta)\), and the overlap identity is
\begin{equation}
 \chi_\gamma\circ h_j=h_j^{S^1}\circ\overline\chi_j.
 \label{cusp-spec-finite-overlap-identity}
\end{equation}

The diagonal action on
\((J'_j)^{\times_{\Delta_j}4}\) is free, and the same coprimality argument
proves that \(C_j^{[4]}\) is smooth.  Every fibre of the resulting map is
a \(\mathbb T^{12}\)-torsor; over \(D_j^*\) its structure torus is the
fourfold fibre product of \(\mathcal H\).  No group zero is asserted at
the affine filling.
In local etale coordinates along the central fibre, the ordinary filled
fibre product has equations
\begin{equation}
 s_i^{m_j}=s_1^{m_j},\qquad i=2,3,4.
 \label{cusp-spec-finite-fibre-product-branches}
\end{equation}
It has \(m_j^3\) smooth branches along the reduced central fibre.  The
simultaneous-root model separates them, is smooth and hence normal, and is
an isomorphism away from the central fibre.  The map \(\nu_j\) is proper:
both spaces are proper over \(D_j\), and the inverse image of a compact
subset of the target is closed in the inverse image of its compact image
in \(D_j\).  It is quasi-finite, since its fibres are singletons away from
the central fibre and have at most \(m_j^3\) points over the central
fibre.  Hence \(\nu_j\) is finite.  The universal property of
normalization now proves that it is the normalization of the reduced
ordinary fibre product.  On the central fibre the relative root choices form
\((\mathbb Z/m_j)^4/\operatorname{diag}(\mathbb Z/m_j)\), of order
\(m_j^3\), which proves the asserted degrees.

On \(D_j^*\), the source identification is \((h_j)^{\times4}\), and
\cref{cusp-spec-finite-overlap-identity} identifies the target overlap
with \((h_j^{S^1})^{\times4}\).  These are precisely the fourfolds of
the one-copy overlap maps and therefore satisfy the cocycle identities.
They glue the two local normalizations and Seifert targets to
\(J^{[4]}\) and to
\(B^\circ\times(\mathbb R/\mathbb Z)^4\), respectively.  Equivalently,
the source gluing follows from uniqueness of normalization.

The fixed product-target trivialization does not extend over \(X\), even if
base winding is allowed.  Indeed, in the presentation of
\cref{prop:geometric-pi1}, let a homomorphism to \(\mathbb Z\) restrict to
\(n\varepsilon_\gamma\) on \(\Lambda\) and take values \(a,b\) on
\(\widehat\rho_1,\widehat\rho_2\).  The three attachment relations give
\[
 3a=n,\qquad4b=-n,\qquad a+b=0.
\]
The first two give \(a=n/3\), \(b=-n/4\), while the third gives
\(n/12=0\).  Hence \(n=0\).  This is the exact global obstruction behind
the fact that the two Seifert targets do not combine into a fixed-circle
map on \(X\).

\medskip\noindent\emph{The explicit toric cusp obstruction.}\par
On the positive radial ray, write
\[
 Z_T=iTB_0+C(t_c),\qquad
 B_0=\begin{pmatrix}0&1\\-1&0\end{pmatrix},\qquad
 t_c=e^{-2\pi T}.
\]
For \(a=(a_1,a_2)^{\mathsf T}\in\mathbb R^2\), take the point of the
regular fibre represented by \(Z_Ta\), with zero short-period coordinates,
and exponentiate it to the dense torus of the cusp model:
\[
 x_T(a)=\exp(2\pi iZ_Ta)\in(\mathbb C^*)^2.
\]
Its logarithmic position is
\begin{equation}
 y_T(a)=
 \frac{(\log|x_{T,1}(a)|,\log|x_{T,2}(a)|)}{\log|t_c|}
 =B_0a+\frac{\operatorname{Im}C(t_c)a}{T}.
 \label{cusp-spec-cusp-nonextension-position}
\end{equation}
The circle character has the constant value
\begin{equation}
 \chi_{\gamma,T}(x_T(a))=a_1\bmod\mathbb Z.
 \label{cusp-spec-cusp-nonextension-character}
\end{equation}

Use the affine toric chart of
\(\sigma=\operatorname{conv}(0,e_1,e_2)\).  Its coordinates are
\begin{equation}
 z_0=\frac{t_c}{x_1x_2},\qquad z_1=x_1,\qquad z_2=x_2.
 \label{cusp-spec-cusp-nonextension-chart}
\end{equation}
If \(B_0a=(y_1,y_2)\) is in the interior of \(\sigma\), then
\(y_1>0\), \(y_2>0\), and \(1-y_1-y_2>0\).  For all sufficiently large
\(T\), the three absolute values in
\cref{cusp-spec-cusp-nonextension-chart} are
\[
 |t_c|^{1-y_{T,1}-y_{T,2}},\qquad
 |t_c|^{y_{T,1}},\qquad |t_c|^{y_{T,2}},
\]
and all tend to zero.  Every such path converges to the same toric fixed
point and hence to the same triple point of \(W\).

Take
\[
 a=\left(-\frac14,\frac14\right),\qquad
 a'=\left(-\frac12,\frac14\right).
\]
Then
\[
 B_0a=\left(\frac14,\frac14\right),\qquad
 B_0a'=\left(\frac14,\frac12\right)
\]
are both in the interior of \(\sigma\).  The paths converge to the same
point of \(W\), whereas
\[
 \chi_{\gamma,T}(x_T(a))=\frac34,\qquad
 \chi_{\gamma,T}(x_T(a'))=\frac12.
\]
No continuous extension of the literal map \(\chi_\gamma\) exists.
Choosing identical marked radial paths in the other three factors and
using these two paths in one factor proves the same assertion for
\(\chi_\gamma^{[4]}\).

The radial family of translations by \(\widehat u/2\) in \(\mathcal H\)
fails to extend as well.  Take two paths
whose limiting logarithmic positions are
\[
 y=\left(\frac18,\frac18\right),\qquad
 y'=\left(\frac38,\frac38\right).
\]
Both approach the fixed point of
\(\operatorname{conv}(0,e_1,e_2)\).  Translation by
\(\widehat u/2\in\mathcal H_T\) adds the exact displacement
\[
 \frac12B_0e_2+
 \frac1{2T}\operatorname{Im}C(t_c)e_2
 \longrightarrow\left(\frac12,0\right)
\]
to logarithmic position.  The translated limits are
\[
 \left(\frac58,\frac18\right)
 \quad\hbox{and}\quad
 \left(\frac78,\frac38\right).
\]
The first lies in the interior of
\(\operatorname{conv}(0,e_1,e_2)\), while the second lies in the interior
of \(\operatorname{conv}(e_1,e_2,e_1+e_2)\).  These two triangle types are
the two translation orbits in the \(A_2\) triangulation and give the two
distinct triple points of \(W\), as proved in
\cref{thm:candidate-cusp-geometry}.  Thus one regular-fibre translation
sends two paths with the same limit to paths with different limits; it has
no continuous extension to \(N_0\).

This pointwise obstruction must not be confused with the cohomology class.
By \cref{prop:geometric-pi1},
\[
 \pi_1(N_0)=\Lambda/K_\Lambda,
\]
and \(\varepsilon_\gamma\) kills \(K_\Lambda\).  It therefore defines an
integral \(H^1\)-class on \(N_0\), hence a map to the Eilenberg--Mac Lane
space \(S^1\) whose restriction is homotopic to \(\chi_\gamma\).  The
explicit paths prove that no such representative agrees with
\(\chi_\gamma\) on the whole punctured cusp.  The phase torus
\((K_\Lambda\otimes\mathbb R)/K_\Lambda\) extends as the compact real torus
action of the toric model, generally with boundary stabilizers; the full
radial \(\widehat u/2\)-translation and the full quotient do not.

\medskip\noindent\emph{Products and smash quotients.}\par
The integral Kunneth theorem gives
\[
 H_k(X^n;\mathbb Z)\cong
 \begin{cases}
  \mathbb Z^{\binom nj},&k=6j,\quad0\le j\le n,\\
  0,&\text{otherwise},
 \end{cases}
 \qquad \chi(X^n)=2^n,
\]
which proves
\cref{cusp-spec-product-not-sphere,cusp-spec-four-product-not-sphere}.
The synchronized regular \(n\)-fold fibre product has dimension
\[
 6n-2(n-1)=4n+2,
\]
which gives \(18\) for four copies.

A product map factors through a smash quotient if and only if it is
constant on the corresponding fat wedge.  On
\(F_T\times\{0\}\), the local two-copy quotient is
\[
 x\longmapsto(\pi_T(x),0),
\]
which is nonconstant.  Thus the literal conclusion is that
\(\pi_T\times\pi_T\) does not factor through \(F_T\wedge F_T\); it would
be a type error to call it a map on \(X\wedge X\).  The restriction of
\(\chi_{\gamma,T}^{[4]}\) to \(F_T\times\{0\}^3\) is similarly
\[
 x\longmapsto(\chi_{\gamma,T}(x),0,0,0),
\]
so it does not factor through \(F_T^{\wedge4}\).

There is no global extension which repairs the domain mismatch:
\(H^1(X^n;\mathbb Z)=0\), whereas both local quotient maps pull nonzero
target \(H^1\)-classes back to nonzero classes on the cusp fibre.  Hence no
map on \(X^2\) or \(X^4\) has the required restriction, and precomposition
with the smash collapses in \cref{cusp-spec-smash-spheres} cannot produce
the local quotient maps.

Identify the metric quotient circle with
\(C_T=\mathbb R/\ell_T\mathbb Z\) by
\(\theta\mapsto\ell_T\theta\).  For
\cref{cusp-spec-target-smash-map}, target-smash collapse is constant
whenever one source coordinate is the identity, so the modified map does
factor through \(F_T^{\wedge4}\).  Its chosen base point has distance
\(\ell_T/2\) from the collapsed target fat wedge.  If two points lie
within distance \(R\) of it, a route through the collapsed point has
length at least \(\ell_T-2R\), while their direct Euclidean distance is at
most \(2R\).  When \(\ell_T>4R\), neither the collapse nor circle
wrap-around shortens distances in that ball.

Let \(u_i\in H^1(F_T^4;\mathbb Z)\) be the pullback of the generator from
the \(i\)-th quotient circle.  The orientation class of
\(S^4=C_T^{\wedge4}\) pulls back under
\cref{cusp-spec-target-smash-map} to
\[
 \pm u_1\smile u_2\smile u_3\smile u_4\ne0.
\]
But \(H^4(X^4;\mathbb Z)=0\).  Thus even this modified \(S^4\)-target map
does not extend to \(X^4\), and a fortiori does not arise from a map on
\(X^{\wedge4}\cong S^{24}\).

For completeness, Borel--Serre prove that \(S^{2n}\) admits no
almost-complex structure for \(n\ge4\).  Their argument applies to every
differentiable structure on the sphere: on \(S^{2n}\), all lower Chern
classes vanish, while their reduced-power calculation forces the top
Chern class to vanish modulo a suitable odd prime \(p<n\); on the other
hand that class is the Euler class, namely \(2h\) for a generator
\(h\in H^{2n}(S^{2n};\mathbb Z)\), a contradiction.  Taking \(n=6\) and
\(n=12\) gives the asserted obstructions for \(S^{12}\) and \(S^{24}\).
This topological obstruction does not itself prohibit the real
four-circle target \(G_T\), the target-smash \(S^4\), or their metric
limits.

Finally,
\(X*X\cong S^{13}\) and \(X^{*4}\cong S^{27}\), so joins do not produce
the proposed \(S^{12}\) or \(S^{24}\) dimensions.  Suspensions can produce
spheres of those dimensions only by discarding the torus-fibration maps.
\end{proof}
